% =============================================================================
%  PREFÁCIO — RELAÇÃO COM O VOLUME 1 E JORNADA DE APRENDIZADO
% =============================================================================
\newpage
\chapter*{Prefácio}
\addcontentsline{toc}{chapter}{Prefácio}

\noindent Esta obra constitui o \textbf{Volume 2} da \textbf{Coleção Tecnologia na Saúde}, dando continuidade direta aos fundamentos estabelecidos no \textbf{Volume 1 --- ``Gêmeos Digitais na Odontologia: Revolução, Impacto Social e o Papel dos Ecossistemas Abertos de IA''}.

\vspace{0.5cm}

\noindent Enquanto o Volume 1 ofereceu uma visão panorâmica e estrutural --- abrangendo desde a definição conceitual dos gêmeos digitais até o arcabouço regulatório e o impacto social ---, o presente volume assume uma premissa radicalmente diferente: \textbf{o leitor não precisa mais se perguntar ``por que'' construir gêmeos digitais, mas sim ``como'' fazê-lo com as ferramentas que já existem, hoje, gratuitamente.}

\vspace{0.5cm}

\noindent Esta mudança de paradigma foi impulsionada por uma constatação empírica: entre 2024 e 2026, o ecossistema open-source amadureceu a ponto de disponibilizar \textbf{25 ou mais ferramentas clinicamente validadas} para cada etapa da construção de gêmeos digitais periodontais. Do DentalSegmentator (segmentação CBCT com nnU-Net, 470 exames de treino) ao Periomod (modelagem preditiva periodontal com K-Fold), do FHIR Brasil (interoperabilidade com a RNDS) ao CaTGO (inferência causal com Transformers) --- o gargalo não é mais tecnológico, mas de \textbf{metodologia e integração}.

\vspace{0.5cm}

\noindent O \textbf{Framework SUS-Twin} surge como resposta a este gargalo: uma metodologia em 6 camadas que organiza e integra estas ferramentas em um pipeline coerente, validado por 312 testes TDD (100\% de aprovação) e alinhado às diretrizes TRIPOD-AI, CLAIM-AI e RDC 657/2022 da ANVISA.

\vspace{0.5cm}

\noindent O Volume 2 está organizado como uma \textbf{jornada progressiva de aprendizado} em quatro níveis:

\begin{itemize}
    \item \textbf{Nível 0 --- Vibe Coding (Caps. 1--4)}: Sem programação. Usando 3D Slicer, DentalSegmentator e IA generativa, o leitor visualiza seu primeiro gêmeo digital em 30 minutos.
    \item \textbf{Nível 1 --- Implementação Guiada (Caps. 5--9)}: Modificação de scripts Python existentes, pipelines de segmentação com MONAI e modelagem com Periomod.
    \item \textbf{Nível 2 --- Pesquisa Reprodutível (Caps. 10--13)}: Projeto de experimentos, validação temporal com K-Fold, cálculo de SROI e publicação seguindo TRIPOD-AI.
    \item \textbf{Nível 3 --- PhD e Inovação (Caps. 14--17)}: Adaptação do framework CaTGO, contribuição com código open-source e publicação em periódicos Qualis A1.
\end{itemize}

\vspace{0.5cm}

\noindent Cada capítulo contém projetos práticos com código executável, métricas de validação e saída esperada para autoavaliação. O leitor pode começar de onde estiver.

\vspace{0.5cm}

\noindent Esta estrutura em dois volumes reflete a maturação do ecossistema de pesquisa: o Volume 1 estabelece o ``o quê'' e o ``porquê''; o Volume 2 responde ao ``como'', oferecendo código, protocolos e validação empírica para a implementação concreta no contexto da saúde pública brasileira.

\vspace{1cm}

\begin{flushright}
\textit{Marcelo Claro Laranjeira}\\
Arquiteto-Chefe, OpenCode Ecosystem\\
Crateús, Ceará, 2026
\end{flushright}
